\documentclass[12pt,a4paper]{article}

% ─── Packages ───────────────────────────────────────────────────────────────
\usepackage[margin=2.5cm]{geometry}
\usepackage{amsmath,amssymb,amsthm}
\usepackage{mathtools}
\usepackage{bm}
\usepackage{booktabs}
\usepackage{array}
\usepackage{xcolor}
\usepackage{hyperref}
\usepackage{enumitem}
\usepackage{titlesec}
\usepackage{parskip}
% bbm not available in tectonic; use \boldsymbol{1} for indicator

% ─── Formatting ─────────────────────────────────────────────────────────────
\hypersetup{colorlinks=true, linkcolor=blue!70!black, urlcolor=blue!70!black}
\titleformat{\section}{\large\bfseries}{}{0em}{}[\titlerule]
\titleformat{\subsection}{\normalsize\bfseries}{\thesubsection}{0.7em}{}

\newcommand{\R}{\mathbb{R}}
\newcommand{\bx}{\bm{x}}
\newcommand{\bw}{\bm{w}}
\newcommand{\by}{\bm{y}}
\newcommand{\ba}{\bm{\alpha}}
\newcommand{\ind}[1]{\boldsymbol{1}\!\left[#1\right]}
\newcommand{\norm}[1]{\left\lVert #1 \right\rVert}

\DeclareMathOperator*{\argmax}{arg\,max}
\DeclareMathOperator*{\argmin}{arg\,min}

% ─── Document ────────────────────────────────────────────────────────────────
\begin{document}

\begin{center}
  {\LARGE\bfseries AI6012 Machine Learning --- Assignment Solution}\\[0.4em]
  {\large Total: 25 marks}\\[0.2em]
  {\small\color{gray}Compiled \today}
\end{center}

\vspace{1em}
\hrule
\vspace{1em}

% ═══════════════════════════════════════════════════════════════════════════
\section{Question 1 (10 marks): Multi-class Regularised Logistic Regression}
% ═══════════════════════════════════════════════════════════════════════════

\subsection*{Setup}

Consider a $C$-class classification problem with training set $\{(\bx_i, y_i)\}_{i=1}^N$
where $\bx_i \in \R^d$ and $y_i \in \{1, \ldots, C\}$.
The weight matrix is $W = [\bw^{(1)}, \ldots, \bw^{(C)}]$,
$\bw^{(c)} \in \R^d$.
The softmax probabilistic model is:
\begin{equation}
  P(y = c \mid \bx;\, W)
  \;=\;
  \frac{\exp\!\bigl((\bw^{(c)})^\top \bx\bigr)}
       {\displaystyle\sum_{j=1}^{C} \exp\!\bigl((\bw^{(j)})^\top \bx\bigr)},
  \qquad c = 1, \ldots, C.
  \label{eq:softmax}
\end{equation}

% ──────────────────────────────────────────────────────────────────────────
\subsection{Part 1: Regularised Negative Log-Likelihood Loss}
% ──────────────────────────────────────────────────────────────────────────

The negative log-likelihood (cross-entropy) over the dataset, augmented with an
$\ell_2$ regularisation penalty, is:
\begin{equation}
  \mathcal{L}(W)
  \;=\;
  -\sum_{i=1}^{N} \log P(y_i \mid \bx_i;\, W)
  \;+\;
  \frac{\lambda}{2}\sum_{c=1}^{C}\norm{\bw^{(c)}}_{2}^{2}.
  \label{eq:loss_raw}
\end{equation}

Using the one-hot indicator notation $\ind{y_i = c}$ (equals 1 iff $y_i = c$,
else 0) and expanding the log-probability:
\begin{align}
  \log P(y_i \mid \bx_i;\, W)
  &= \sum_{c=1}^{C} \ind{y_i = c}
     \Bigl[(\bw^{(c)})^\top \bx_i
           - \log\textstyle\sum_{j=1}^{C}\exp\!\bigl((\bw^{(j)})^\top \bx_i\bigr)
     \Bigr].
\end{align}

Substituting into \eqref{eq:loss_raw}:
\begin{equation}
  \boxed{
  \mathcal{L}(W)
  \;=\;
  -\sum_{i=1}^{N}\sum_{c=1}^{C}
     \ind{y_i = c}
     \Bigl[(\bw^{(c)})^\top \bx_i
           - \log\sum_{j=1}^{C}\exp\!\bigl((\bw^{(j)})^\top \bx_i\bigr)
     \Bigr]
  \;+\;
  \frac{\lambda}{2}\sum_{c=1}^{C}\norm{\bw^{(c)}}_{2}^{2}.
  }
  \label{eq:loss_full}
\end{equation}

% ──────────────────────────────────────────────────────────────────────────
\subsection{Part 2: Gradient w.r.t.\ \texorpdfstring{$\bw^{(c)}$}{w(c)}}
% ──────────────────────────────────────────────────────────────────────────

\paragraph{Notation.}
Define for sample $i$ the pre-activation $a_k^{(i)} = (\bw^{(k)})^\top \bx_i$
and the softmax probability
$p_c^{(i)} = P(y = c \mid \bx_i;\, W)$.

\paragraph{Step 1: Isolate the terms that depend on $\bw^{(c)}$.}

The regularisation contribution gives immediately $\lambda \bw^{(c)}$.
The data term involves every $\log P(y_i \mid \bx_i;\, W)$; we differentiate it
w.r.t.\ $\bw^{(c)}$:
\begin{align}
  \frac{\partial}{\partial \bw^{(c)}}
  \log P(y_i \mid \bx_i;\, W)
  &= \sum_{k=1}^{C} \ind{y_i = k}
     \frac{\partial}{\partial \bw^{(c)}}
     \Bigl[a_k^{(i)} - \log\sum_{j=1}^{C}e^{a_j^{(i)}}\Bigr].
  \label{eq:grad_start}
\end{align}

\paragraph{Step 2: Differentiate each component.}

\begin{itemize}[leftmargin=*]
  \item \textbf{Linear term:}
    $\dfrac{\partial a_k^{(i)}}{\partial \bw^{(c)}}
     = \ind{k = c}\,\bx_i$,
    since $a_k^{(i)} = (\bw^{(k)})^\top\bx_i$ only depends on $\bw^{(c)}$
    when $k = c$.

  \item \textbf{Log-sum-exp term:}
    \begin{align}
      \frac{\partial}{\partial \bw^{(c)}} \log\sum_{j=1}^{C}e^{a_j^{(i)}}
      &= \frac{1}{\sum_{j}e^{a_j^{(i)}}}
         \cdot\frac{\partial}{\partial \bw^{(c)}}
         \sum_{j=1}^{C}e^{a_j^{(i)}} \notag\\
      &= \frac{e^{a_c^{(i)}}}{\sum_{j}e^{a_j^{(i)}}}
         \cdot\frac{\partial a_c^{(i)}}{\partial \bw^{(c)}} \notag\\
      &= p_c^{(i)}\,\bx_i.
    \end{align}
\end{itemize}

\paragraph{Step 3: Combine.}

Substituting back into \eqref{eq:grad_start}:
\begin{align}
  \frac{\partial}{\partial \bw^{(c)}} \log P(y_i \mid \bx_i;\, W)
  &= \sum_{k=1}^{C} \ind{y_i = k}
     \bigl(\ind{k = c} - p_c^{(i)}\bigr)\bx_i \notag\\
  &= \Bigl[\ind{y_i = c}
     \;-\; p_c^{(i)}\underbrace{\sum_{k=1}^{C}\ind{y_i = k}}_{=1}\Bigr]\bx_i \notag\\
  &= \bigl(\ind{y_i = c} - p_c^{(i)}\bigr)\bx_i.
\end{align}

\paragraph{Step 4: Full gradient.}

Taking the negative and adding the regularisation term:
\begin{equation}
  \boxed{
  \frac{\partial \mathcal{L}}{\partial \bw^{(c)}}
  \;=\;
  \sum_{i=1}^{N}\bigl(p_c^{(i)} - \ind{y_i = c}\bigr)\bx_i
  \;+\;
  \lambda\,\bw^{(c)},
  \qquad c = 1, \ldots, C.
  }
  \label{eq:gradient}
\end{equation}

\textit{Interpretation:} Each term is the product of the prediction
error $(p_c^{(i)} - \ind{y_i=c})$ and the feature vector $\bx_i$,
summed over all samples, plus the weight-decay term $\lambda\bw^{(c)}$.

% ──────────────────────────────────────────────────────────────────────────
\subsection{Part 3: Gradient Descent Update Rules}
% ──────────────────────────────────────────────────────────────────────────

With learning rate $\eta > 0$, the update at step $(t+1)$ is:
\begin{equation}
  \bw^{(c)}_{(t+1)}
  \;=\;
  \bw^{(c)}_{(t)}
  \;-\;
  \eta\,\frac{\partial \mathcal{L}}{\partial \bw^{(c)}}\Bigg|_{(t)}.
\end{equation}

Substituting \eqref{eq:gradient}:
\begin{equation}
  \boxed{
  \bw^{(c)}_{(t+1)}
  \;=\;
  (1 - \eta\lambda)\,\bw^{(c)}_{(t)}
  \;-\;
  \eta\sum_{i=1}^{N}
     \Bigl(p_{c,{(t)}}^{(i)} - \ind{y_i = c}\Bigr)\bx_i,
  \qquad c = 1, \ldots, C,
  }
  \label{eq:update}
\end{equation}
where $p_{c,(t)}^{(i)} = P(y=c \mid \bx_i;\, W_{(t)})$ is the predicted
probability computed at iteration $t$.
The factor $(1-\eta\lambda)$ is the standard \emph{weight-decay} shrinkage
introduced by the $\ell_2$ regulariser.

% ═══════════════════════════════════════════════════════════════════════════
\newpage
\section{Question 2 (5 marks): SVM Experiments with scikit-learn on a9a}
% ═══════════════════════════════════════════════════════════════════════════

\subsection*{Experimental Setup}

\begin{itemize}[leftmargin=*]
  \item \textbf{Dataset}: LIBSVM \texttt{a9a}
        (train: 32,561 samples, test: 16,281 samples; 123 features).
  \item \textbf{Cross-validation}: 3-fold CV
        (\texttt{n\_splits=3}, \texttt{shuffle=True}, \texttt{random\_state=42}).
  \item \textbf{Implementation}: \texttt{sklearn.svm.SVC} (\texttt{scikit-learn}).
  \item \textbf{Metric}: Mean validation accuracy over 3 folds.
\end{itemize}

% ──────────────────────────────────────────────────────────────────────────
\subsection{Table 1: Linear Kernel --- Mean 3-fold Validation Accuracy}
% ──────────────────────────────────────────────────────────────────────────

\begin{table}[h!]
\centering
\begin{tabular}{cc}
\toprule
$C$ & Mean Validation Accuracy \\
\midrule
0.01 & 0.844784 \\
0.05 & 0.847179 \\
\textbf{0.10} & \textbf{0.847978} \\
0.50 & 0.847947 \\
1.00 & 0.847855 \\
\bottomrule
\end{tabular}
\caption{Linear kernel 3-fold cross-validation results.
         Best: $C^* = 0.1$ with mean accuracy $0.847978$.}
\label{tab:linear}
\end{table}

% ──────────────────────────────────────────────────────────────────────────
\subsection{Table 2: RBF Kernel --- Mean 3-fold Validation Accuracy (5$\times$5 Grid)}
% ──────────────────────────────────────────────────────────────────────────

\begin{table}[h!]
\centering
\setlength{\tabcolsep}{6pt}
\begin{tabular}{c|ccccc}
\toprule
$C \;\backslash\; \gamma$ & 0.01 & 0.05 & 0.10 & 0.50 & 1.00 \\
\midrule
0.01 & 0.759190 & 0.820061 & 0.819508 & 0.759190 & 0.759190 \\
0.05 & 0.831025 & 0.835601 & 0.834157 & 0.788858 & 0.759190 \\
0.10 & 0.837904 & 0.839655 & 0.839041 & 0.805504 & 0.761893 \\
0.50 & 0.843555 & 0.844753 & 0.846411 & 0.832775 & 0.788950 \\
1.00 & 0.845060 & 0.846964 & \textbf{0.847241} & 0.834679 & 0.797549 \\
\bottomrule
\end{tabular}
\caption{RBF kernel 3-fold cross-validation results.
         Best: $C^* = 1.0,\; \gamma^* = 0.1$ with mean accuracy $0.847241$.}
\label{tab:rbf}
\end{table}

% ──────────────────────────────────────────────────────────────────────────
\subsection{Table 3: Test Set Evaluation}
% ──────────────────────────────────────────────────────────────────────────

\paragraph{Model selection.}
Comparing the best configurations:
\begin{itemize}[leftmargin=*]
  \item Linear, $C=0.1$: mean val.\ accuracy $= 0.847978$.
  \item RBF, $C=1.0, \gamma=0.1$: mean val.\ accuracy $= 0.847241$.
\end{itemize}
The \textbf{linear kernel with $C=0.1$} is selected (marginally higher
validation accuracy).  It is retrained on the \emph{entire} training set
and evaluated on the held-out test set.

\begin{table}[h!]
\centering
\begin{tabular}{lc}
\toprule
 & Linear kernel, $C = 0.1$ \\
\midrule
Accuracy of SVMs & \textbf{0.850132} \\
\bottomrule
\end{tabular}
\caption{Test results of SVC on the a9a test set (in accuracy).
         The linear kernel with $C=0.1$ was retrained on all 32,561 training
         samples and evaluated on the held-out test set of 16,281 samples.}
\label{tab:test}
\end{table}

\paragraph{Observation.}
The linear kernel achieves \textbf{85.01\%} test accuracy.
This is consistent with the finding that the \texttt{a9a} dataset is nearly
linearly separable in the original 123-dimensional feature space; the RBF
kernel offers negligible benefit and adds the extra hyperparameter $\gamma$.

% ═══════════════════════════════════════════════════════════════════════════
\newpage
\section{Question 3 (5 marks): Soft-Margin SVM as Structural Risk Minimisation}
% ═══════════════════════════════════════════════════════════════════════════

\subsection*{Primal Soft-Margin SVM}

The standard soft-margin SVM solves:
\begin{equation}
  \min_{\bw,\, b,\, \bm{\xi}}\;
  \frac{1}{2}\norm{\bw}^2 + C\sum_{i=1}^{N}\xi_i
  \label{eq:svm_primal}
\end{equation}
\begin{equation*}
  \text{subject to}\quad
  y_i(\bw^\top\bx_i + b) \;\geq\; 1 - \xi_i,
  \qquad
  \xi_i \;\geq\; 0,
  \qquad i = 1,\ldots,N.
\end{equation*}

Let $f(\bx) = \bw^\top\bx + b$.

\subsection*{Hinge Loss Definition}

\begin{equation}
  \boxed{
  L_{\text{hinge}}(y,\, f(\bx))
  \;=\;
  \max\!\bigl(0,\; 1 - y\cdot f(\bx)\bigr).
  }
  \label{eq:hinge}
\end{equation}

This penalises predictions that are either on the wrong side of the margin
($y f(\bx) < 1$) linearly, and gives zero loss to correct predictions with
sufficient margin ($y f(\bx) \geq 1$).

\subsection*{Algebraic Equivalence}

\paragraph{Step 1: Express $\xi_i$ in closed form.}

From the two constraints, for each sample $i$:
\begin{equation*}
  \xi_i \;\geq\; 1 - y_i f(\bx_i)
  \qquad\text{and}\qquad
  \xi_i \;\geq\; 0.
\end{equation*}
Combining: $\xi_i \geq \max(0,\; 1 - y_i f(\bx_i))$.

\paragraph{Step 2: Identify the optimal slack.}

Since $C > 0$, the objective is minimised by making each $\xi_i$ as small as
possible.  Hence at optimality:
\begin{equation}
  \xi_i^* \;=\; \max\!\bigl(0,\; 1 - y_i f(\bx_i)\bigr)
            \;=\; L_{\text{hinge}}\!\bigl(y_i,\, f(\bx_i)\bigr).
  \label{eq:xi_opt}
\end{equation}

\paragraph{Step 3: Substitute back into the objective.}

Replacing $\xi_i$ with $\xi_i^*$ in \eqref{eq:svm_primal} and dropping the
(now redundant) slack variable constraints:
\begin{equation}
  \min_{\bw,\, b}
  \;\frac{1}{2}\norm{\bw}^2
  \;+\; C\sum_{i=1}^{N}\max\!\bigl(0,\; 1 - y_i f(\bx_i)\bigr).
  \label{eq:erm_form}
\end{equation}

\paragraph{Step 4: SRM formulation.}

Rearranging \eqref{eq:erm_form} into the Empirical Structural Risk Minimisation
(SRM) form:
\begin{equation}
  \boxed{
  \min_{\bw,\, b}
  \underbrace{C\sum_{i=1}^{N}
     L_{\text{hinge}}\!\bigl(y_i,\, f(\bx_i)\bigr)}_{\text{Empirical Risk}}
  \;+\;
  \underbrace{\frac{1}{2}\norm{\bw}^2}_{\text{Regularisation (Complexity Penalty)}}.
  }
  \label{eq:srm}
\end{equation}

The two problems \eqref{eq:svm_primal} and \eqref{eq:srm} are
\textbf{exactly equivalent}: given any feasible point $(\bw, b)$, the objectives
coincide because the slack variables are uniquely determined by \eqref{eq:xi_opt}.
The trade-off parameter $C$ controls the balance between minimising training
hinge loss (empirical risk) and controlling the weight norm (structural
complexity of the model class).

% ═══════════════════════════════════════════════════════════════════════════
\newpage
\section{Question 4 (5 marks): Kernelised Ridge Regression}
% ═══════════════════════════════════════════════════════════════════════════

\subsection*{Setup and Kernel Extension Idea}

To solve \emph{nonlinear} regression, we replace every input $\bx_i \in \R^d$
with its image $\phi(\bx_i)$ under a (possibly infinite-dimensional) feature
map $\phi: \R^d \to \mathcal{H}$.  All inner products then become kernel
evaluations:
\[
  \phi(\bx_i)^\top \phi(\bx_j) \;=\; k(\bx_i,\, \bx_j).
\]
We now derive the regularised solution entirely in terms of $k$, so $\phi$ need
never be computed explicitly.

Let $\Phi \in \R^{N\times D}$ (or infinite) be the mapped data matrix whose
$i$-th row is $\phi(\bx_i)^\top$.  The regularised regression problem becomes:
\begin{equation}
  \min_{\bw}\;
  \norm{\Phi\bw - \by}^2 + \lambda\norm{\bw}^2,
  \label{eq:ridge_primal}
\end{equation}
with closed-form \textbf{primal solution}:
\begin{equation}
  \bw = (\Phi^\top \Phi + \lambda I)^{-1} \Phi^\top \by.
  \label{eq:primal_sol}
\end{equation}
(For the linear case $\phi(\bx)=\bx$, $\Phi = X$ and this recovers the
standard ridge solution on L3 p.\,50.)
We now apply the \textbf{kernel trick} by rewriting this solution so that
only the $N\times N$ Gram matrix $K_{ij}=k(\bx_i,\bx_j)$ appears.

% ──────────────────────────────────────────────────────────────────────────
\subsection{Part 1: Representer Theorem --- \texorpdfstring{$\bw = \Phi^\top \ba$}{w = Phi' alpha}}
% ──────────────────────────────────────────────────────────────────────────

The gradient (first-order optimality condition) of \eqref{eq:ridge_primal} is:
\begin{equation}
  2\Phi^\top(\Phi\bw - \by) + 2\lambda\bw = \bm{0}.
\end{equation}
Rearranging:
\begin{align}
  \lambda\bw
  &= \Phi^\top\by - \Phi^\top\Phi\bw
   = \Phi^\top\!\bigl(\by - \Phi\bw\bigr). \notag
\end{align}
Dividing both sides by $\lambda$:
\begin{equation}
  \bw \;=\; \Phi^\top\underbrace{\frac{1}{\lambda}\bigl(\by - \Phi\bw\bigr)}_{\displaystyle\ba}.
  \label{eq:representer}
\end{equation}

Therefore $\bw = \Phi^\top\ba$, where $\ba \in \R^N$ is a vector of
\textbf{dual coefficients} (one per training sample).  This is the
\textbf{Representer Theorem}: the optimal solution $\bw$ lies in the span
of the mapped training inputs $\{\phi(\bx_i)\}_{i=1}^N$.

% ──────────────────────────────────────────────────────────────────────────
\subsection{Part 2: Closed-form Solution for \texorpdfstring{$\ba$}{alpha}}
% ──────────────────────────────────────────────────────────────────────────

Substitute $\bw = \Phi^\top\ba$ into $\ba = \frac{1}{\lambda}(\by - \Phi\bw)$:
\begin{align}
  \ba
  &= \frac{1}{\lambda}\bigl(\by - \Phi\Phi^\top\ba\bigr)
   = \frac{1}{\lambda}\bigl(\by - K\ba\bigr),
\end{align}
where $K = \Phi\Phi^\top \in \R^{N\times N}$ is the \textbf{kernel (Gram) matrix}
with entries
\[
  K_{ij} = \phi(\bx_i)^\top\phi(\bx_j) = k(\bx_i, \bx_j).
\]
Multiplying both sides by $\lambda$:
\begin{equation}
  \lambda\ba = \by - K\ba
  \;\;\Longrightarrow\;\;
  (K + \lambda I_N)\ba = \by.
\end{equation}
Solving:
\begin{equation}
  \boxed{
  \ba = (K + \lambda I_N)^{-1}\by.
  }
  \label{eq:alpha_sol}
\end{equation}

Since $K$ is symmetric positive semi-definite and $\lambda > 0$, the matrix
$(K + \lambda I_N)$ is strictly positive definite and hence invertible.
This is the \textbf{kernelised closed-form solution} (the dual analogue of
the primal solution \eqref{eq:primal_sol}).

% ──────────────────────────────────────────────────────────────────────────
\subsection{Part 3: Kernelised Prediction}
% ──────────────────────────────────────────────────────────────────────────

For a new test input $\bx^* \in \R^d$ with feature map $\phi(\bx^*)$:
\begin{align}
  f(\bx^*)
  &= \bw^\top\phi(\bx^*)
   = (\Phi^\top\ba)^\top\phi(\bx^*) \notag\\
  &= \ba^\top \Phi\,\phi(\bx^*) \notag\\
  &= \ba^\top \bm{k}(\bx^*),
\end{align}
where
$\bm{k}(\bx^*) = \bigl[k(\bx_1, \bx^*),\, k(\bx_2, \bx^*),\, \ldots,\,
k(\bx_N, \bx^*)\bigr]^\top \in \R^N$
collects all kernel evaluations between $\bx^*$ and the training points
(since the $i$-th entry of $\Phi\,\phi(\bx^*)$ equals $\phi(\bx_i)^\top\phi(\bx^*)=k(\bx_i,\bx^*)$).

Writing the prediction explicitly:
\begin{equation}
  \boxed{
  f(\bx^*) = \sum_{i=1}^{N} \alpha_i\, k(\bx_i, \bx^*),
  }
  \label{eq:prediction}
\end{equation}
with $\ba = (K + \lambda I_N)^{-1}\by$.
The prediction \emph{never} requires computing $\phi$ explicitly; only the
kernel function $k(\cdot,\cdot)$ is needed.  The complete algorithm is:
\begin{enumerate}[leftmargin=*]
  \item Compute $K_{ij} = k(\bx_i, \bx_j)$ for all $N^2$ training pairs.
  \item Solve $\ba = (K + \lambda I_N)^{-1}\by$ (kernel ridge regression).
  \item For each test point $\bx^*$, predict via \eqref{eq:prediction}.
\end{enumerate}

\end{document}
